\subsection{Matched FastAPI Comparison}
Table~\ref{tab:api-model-size} reports the matched Vulkan comparison under the original support labels, alongside raw unsupported selections and unvetted original-compatible kills (mutant failures on suites that pass the reference service before checking requirement support). On vague requirements, scaling from 0.5B to 7B lowers the raw unsupported rate from $49/80$ (61.3\%) to $31/67$ (46.3\%) and raises raw supported coverage from $31/48$ to $36/48$. After intersection filtering, supported coverage moves from $18/48$ (58.1\% of raw supported choices retained) to $20/48$ (55.6\% retained), with zero unsupported selections under the original labels, while qualified detections move from $22/51$ to $21/51$. On precise text, filtered coverage moves from $50/81$ to $47/81$ and qualified detections from $36/51$ to $37/51$.

\begin{table*}[t]
\caption{Matched FastAPI model comparison under Vulkan (three seeds: 18 suites and 51 mutant opportunities per row). Compat.\ kills count mutant detections by nonempty suites passing the original service prior to support qualification; Qual.\ kills additionally require every selected assertion to belong to $G_r$.}
\label{tab:api-model-size}
\centering\small
\begin{tabular}{llrrrrrrr}
\toprule
 & & \multicolumn{4}{c}{Raw LLM ($L$)} & \multicolumn{3}{c}{Filtered Intersection ($F=L\cap R$)}\\
\cmidrule(lr){3-6}\cmidrule(lr){7-9}
Variant & Model & Supported & Unsupported & Orig.\ pass & Compat.\ / Qual.\ kills & Supported & Empty & Qual.\ kills\\
\midrule
Precise & 0.5B & 50/81 & 23/73 & 4/18 & 7/51 / \phantom{0}7/51 & 50/81 & 2/18 & 36/51\\
Precise & 7B   & 47/81 & 14/61 & 7/18 & 18/51 / 18/51 & 47/81 & 0/18 & 37/51\\
\midrule
Vague   & 0.5B & 31/48 & 49/80 & 1/18 & 2/51 / \phantom{0}2/51 & 18/48 & 5/18 & 22/51\\
Vague   & 7B   & 36/48 & 31/67 & 8/18 & 23/51 / \phantom{0}5/51 & 20/48 & 4/18 & 21/51\\
\bottomrule
\end{tabular}
\end{table*}

Crucially, on vague requirements the raw 7B model achieves $8/18$ original-passing suites and $23/51$ compatible mutant kills, up from $1/18$ and $2/51$ at 0.5B, yet $18$ of those $23$ detections come from six suites that assert unstated implementation conventions (such as exact \texttt{http\_200}, \texttt{http\_400}, or \texttt{can\_login} checks) and therefore fail support qualification ($5/51$ qualified kills). A benchmark that grades generated tests solely by reference compatibility and mutant kills would credit those $23/51$ detections, masking specification-oracle drift as model capability increases \cite{konstantinou2024,ding2026}.

The larger model also reduces seed-to-seed variance. Across seeds $(11,23,47)$, raw precise supported counts span $[12,24,14]$ at 0.5B (zero of six requirements unanimous across seeds) versus $[15,15,17]$ at 7B (two of six unanimous), with filtered precise kills stabilizing from $[11,16,9]$ to $[12,12,13]$. Meanwhile, the matched Vulkan 0.5B run reproduces the discrete selection mask of the archived CPU run on only $16/36$ FastAPI calls, whereas an exact CPU rerun with the archived settings matches $36/36$ FastAPI and $144/144$ controlled completions. Changing the execution backend and KV-cache prefix reuse under 4-bit quantization and non-zero temperature ($\tau=0.3$) shifts $55.6\%$ of individual 0.5B masks without altering the qualitative filtering effect, reinforcing methodological warnings on inference-stack sensitivity \cite{sallou2024}.

The aggregate shift from $22/51$ to $21/51$ vague qualified detections combines opposing requirement-level effects (Table~\ref{tab:api-scale-by-requirement}, traced to \href{https://github.com/lythean2508/Evaluating-LLM-Selected-Test-Oracles-for-Underspecified-Natural-Language-Requirements/blob/main/followup/results/api-vague-filtered-by-requirement.csv}{\texttt{api-vague-filtered-by-requirement.csv}}). At 7B, requirement~6 improves from $2/9$ to $5/9$ kills and requirement~34 recovers from $0/9$ (three empty suites at 0.5B) to $3/9$, whereas requirement~4 drops from $6/9$ to $2/9$ because 7B omits the broad \texttt{denied} check in seeds~11 and~23 in favor of exact status checks that the vague policy excludes.

\begin{table}[t]
\caption{Filtered results for vague FastAPI requirements across three seeds under Vulkan. Coverage divides supported selections by available supported candidates. Kills count qualified detections over mutant opportunities.}
\label{tab:api-scale-by-requirement}
\centering\small
\begin{tabular}{lrrrr}
\toprule
 & \multicolumn{2}{c}{Coverage} & \multicolumn{2}{c}{Kills}\\
ID & 0.5B & 7B & 0.5B & 7B\\
\midrule
2  & 2/15 & 1/15 & 2/9 & 1/9\\
3  & 3/6  & 3/6  & 6/6 & 6/6\\
4  & 3/3  & 1/3  & 6/9 & 2/9\\
6  & 4/15 & 8/15 & 2/9 & 5/9\\
13 & 6/6  & 4/6  & 6/9 & 4/9\\
34 & 0/3  & 3/3  & 0/9 & 3/9\\
\midrule
Total & 18/48 & 20/48 & 22/51 & 21/51\\
\bottomrule
\end{tabular}
\end{table}

\subsection{Held-out TodoApp Requirements}
Table~\ref{tab:heldout-status} reports status-oracle outcomes on the four held-out TodoApp pairs. On precise text, 7B raises raw supported coverage from $7/24$ to $13/24$, lowers raw unsupported selections from $42/49$ to $19/32$, and raises filtered coverage from $7/24$ to $10/24$ while cutting empty suites from $6/12$ to $2/12$. On vague text, 0.5B and 7B select $5/12$ and $3/12$ supported opportunities before filtering, but the frozen rules match zero status candidates on all four vague phrasings (e.g., \emph{``returns a validation error''}, \emph{``returns the correct todo''}, \emph{``removes it from the database''}, and \emph{``correctly changes its fields''}), forcing $12/12$ empty filtered suites for both models.

\begin{table}[t]
\caption{Held-out TodoApp status assertions (three seeds: 12 suites per row). Coverage uses available supported candidates. Unsupported uses all filtered selections. Empty uses all filtered suites. A dash marks zero selected assertions.}
\label{tab:heldout-status}
\centering\small
\setlength{\tabcolsep}{4pt}
\begin{tabular}{llrrrr}
\toprule
 & & \multicolumn{2}{c}{Coverage} & \multicolumn{2}{c}{Filtered}\\
Variant & Model & Raw & Filtered & Unsupported & Empty\\
\midrule
Precise & 0.5B & 7/24 & 7/24 & 2/9 & 6/12\\
Precise & 7B & 13/24 & 10/24 & 3/13 & 2/12\\
Vague & 0.5B & 5/12 & 0/12 & -- & 12/12\\
Vague & 7B & 3/12 & 0/12 & -- & 12/12\\
\bottomrule
\end{tabular}
\end{table}

In one deterministic pass on precise TodoApp text, the frozen rules cover $7/8$ supported candidates (missing \texttt{http\_200} on requirement~9 because \emph{``Verify response is 200 OK''} omits the keywords \texttt{status} and \texttt{HTTP}) and introduce one unsupported candidate. Specifically, requirement~6 specifies \texttt{200 OK} for \texttt{DELETE /todos/\{id\}} followed by \texttt{404 Not Found} on a subsequent verification \texttt{GET}. Because the document-level regex selector is step-agnostic, \texttt{404} triggers \texttt{denied}, yielding the mutually contradictory rule suite $\{\texttt{http\_200},\texttt{success},\texttt{denied}\}$ ($2\text{xx}\wedge 4\text{xx}$), which violates original-service compatibility and thus breaks the qualification precondition $Q(R,r)$ of Equation~\ref{eq:bound}. Moreover, both 0.5B (in two seeds) and 7B (in all three seeds, selecting $\{\texttt{http\_200},\texttt{denied}\}$) exhibit the same cross-step scope confusion, so correlated error across $L$ and $R$ preserves an unsatisfiable assertion pair in $2/9$ and $3/13$ filtered selections, respectively.

\subsection{Rule Ceiling and Sensitivity to Support Labels}
Table~\ref{tab:partition} decomposes supported FastAPI opportunities into four disjoint subsets across the archived CPU 0.5B run and the matched Vulkan pair. Under the original vague labels, the frozen rules impose an invariant ceiling of $27/48$ ($56.3\%$) supported opportunities across three seeds. Scaling from 0.5B to 7B under Vulkan shrinks model-missed rule matches ($R$ only) from $9$ to $7$ and joint omissions from $8$ to $5$, but simultaneously increases supported model choices vetoed by the lexical gate ($L$ only) from $13$ to $16$. Regardless of model capacity, intersection filtering cannot recover candidates in $G_r\setminus R(r)$. On vague TodoApp, $R(r)=\emptyset$ prior to label scoring, so the $12/12$ empty suites hold under every possible support annotation.

\begin{table}[t]
\caption{Disjoint partition of supported FastAPI candidate opportunities across three seeds under original labels ($G_r$). $L$ only counts supported model choices vetoed by rules; $R$ only counts rule matches omitted by the model; Ceiling is $|G_r\cap R|$.}
\label{tab:partition}
\centering\small
\setlength{\tabcolsep}{3.0pt}
\begin{tabular}{llrrrrrr}
\toprule
Variant & Run / Model & Both & $L$ only & $R$ only & Neither & Ceiling & $|G_r|$\\
\midrule
Precise & 0.5B (CPU) & 47 & 0 & 34 & 0 & 81 & 81\\
Precise & 0.5B (Vulkan) & 50 & 0 & 31 & 0 & 81 & 81\\
Precise & 7B (Vulkan) & 47 & 0 & 34 & 0 & 81 & 81\\
\midrule
Vague & 0.5B (CPU) & 17 & 11 & 10 & 10 & 27 & 48\\
Vague & 0.5B (Vulkan) & 18 & 13 & 9 & 8 & 27 & 48\\
Vague & 7B (Vulkan) & 20 & 16 & 7 & 5 & 27 & 48\\
\bottomrule
\end{tabular}
\end{table}

To test sensitivity to interpretive policy, we replayed the FastAPI suites across all $2^3=8$ combinations of three stricter annotation rules surfaced during an exploratory AI-assisted review: requiring explicit numeric wording for \texttt{http\_404} on vague not-found text (R3), requiring explicit response-body wording rather than state-update wording for \texttt{PUT} fields (four checks in R6), and requiring explicit token wording for the compound \texttt{can\_login} check (R34). Enabling a switch removes the affected check(s) from $G_r$ while keeping all saved selections and API observations fixed.

Under baseline policy \texttt{000}, filtered vague coverage favors 7B over 0.5B ($20/48$ vs.\ $18/48$). Activating the stricter \texttt{PUT}-field policy reverses this ranking across all four corresponding policies (e.g., $14/36$ vs.\ $16/36$ under \texttt{010}, and $14/33$ vs.\ $16/33$ with $6/20$ vs.\ $2/18$ unsupported selections under \texttt{111}), because 7B more frequently selects \texttt{same\_owner} on R6 where rules also match. Conversely, vague qualified detections remain lower at 7B across all eight policies ($21/51$ vs.\ $22/51$ without \texttt{PUT} restriction; $16/51$ vs.\ $20/51$ with it), and precise qualified detections flip from $37/51$ vs.\ $36/51$ to $28/51$ vs.\ $31/51$ whenever explicit token wording is required (\texttt{001}, \texttt{011}, \texttt{101}, \texttt{111}). Model-scaling conclusions on underspecified requirements are therefore policy-contingent, motivating blind multi-rater adjudication before claiming superiority.
